\section{Countably many compatible partitions}\label{sec:countable}

The separations in \Cref{rounds:wide} may depend on the label.
We now use this freedom to construct the partitions. The analytic
construction is performed at finite stages. Its essential finiteness
property concerns each \emph{whole component}, including all its infinite
product ends; local finiteness merely on compact subsets would not
suffice. The finite-label procedure is stated in \Cref{finite:cut}, and its detailed
proof is in \Cref{sec:components,sec:cut,sec:local,sec:enlargement-global,sec:regularization}.
We adapt that proof using individual positive floors and finite
exceptional data on each whole component.

\subsection{Data and the countable construction}

Let $D$ be the length in \eqref{rounds:original-scale}, and let the
centers $x_i$ and the pairs $(A_i,B_i)$ be those fixed in
Section~\ref{sec:width-map}: a maximal $D$-separated set and its ball
pairs, as in \Cref{rounds:wide}. We assume every $B_i$ is nonempty, since otherwise
$\diam_g M\le4D$. There is one label for each center, and an additional
label $0$ used first for radial confinement. The set of labels is finite
or countable; write it as $I=\{0,1,2,\ldots\}$, stopping the enumeration
if necessary. Choose
\begin{equation}\label{countable:floors}
 p_i^*>0,\qquad \sum_{i\in I}p_i^*=\frac1{20}.
\end{equation}
Use \Cref{rounds:wide} with
\begin{equation}\label{countable:widths}
 L_i=10+2\pi\sqrt{2/p_i^*}\qquad(i\ne0).
\end{equation}
We henceforth use the resulting complete metric $h\ge g$, its locally
finite list of logarithms of the warping functions $p_\alpha$, and its stabilized scalar lower
bound $2/5$.

The initial weighted base measure is globally well defined:
\[
 \omega\,\dd\vol_h,\qquad
 \omega=\exp\Bigl(\sum_\alpha p_\alpha\Bigr)>0.
\]
The sum is locally finite. On a buffered compact set, only finitely many
nonunit factors $e^{p_\alpha}$ occur, so ordinary finite warped-product
formulas compute the stabilized scalar, weighted variations, and
invariant test forms there. These formulas agree on overlaps: adjoining
a unit factor changes none of them. In particular, the initial problem
is a weighted problem on the $n$-dimensional base $M$.

Fix numbers
\[
 0=t_0<t_1<t_2<\cdots<1,\qquad t_r\longrightarrow1,
\]
and use as many stages as there are labels. These are budget parameters,
unrelated to a spatial coordinate. Put
\begin{equation}\label{countable:constants}
 \begin{gathered}
 a_r=.60+.03t_r,\qquad b_r=1.12-.05t_r,
 \qquad \lambda_r=\frac14-\frac{t_r}{8},\\
 \Delta_r=t_{r+1}-t_r,\qquad
 e_r=\frac{\Delta_r}{800},\qquad
 \eta_r=\frac{\Delta_r}{200},\qquad k=399,
 \qquad C=\frac{4k}{k+1}=3.99.
 \end{gathered}
\end{equation}
With $t_r=r/N_*$, the entries shared with \eqref{component:constants}
agree with it, $e_r$ being its $e_0$.
Thus $\lambda_r\ge1/8$, $.60\le a_r<.63$, and
$1.07<b_r\le1.12$. In a stagewise
appendix lemma, its symbol $e_0$ denotes the current $e_r$.

We use the component notation and operation of splitting along regular zero sets introduced in
\Cref{sec:finite-cuts}. The countable invariant below replaces the common
finite-label floor by the individual positive floors $p_j^*$.
At stage $r$, the processed labels are $0,\ldots,r-1$ and the remaining
set is $J_r$. A component with $s$ star entries has base dimension $m=n-s$.
Its continuous collapse map to $M$ is denoted by $q_E$. For each
$j\in J_r$ retain an extended-real continuous function $\mu_j$ with
locally Lipschitz finite truncations, a smooth potential
$\Pi_j\ge p_j^*$, and disjoint prescribed open sets $U_j^+,U_j^-$ with
\begin{equation}\label{countable:function-invariant}
 \overline{U_j^\pm}\subset\operatorname{int}\{\mu_j=\pm\infty\},
 \qquad 2|\grad\mu_j|\le a_r\mu_j^2+\Pi_j
 \quad\hbox{almost everywhere on the finite-value locus}.
\end{equation}
For unprocessed center labels require
$q_E^{-1}(A_j)\subset U_j^+$ and $q_E^{-1}(B_j)\subset U_j^-$.
All data are compatible products in the component ends. With
\begin{equation}\label{countable:required}
 G_j=E\setminus(U_j^+\cup U_j^-),\qquad
 D_E=\{x:\#\{j\in J_r:x\in G_j\}\ge m-1\},
\end{equation}
the scalar requirement is
\begin{equation}\label{countable:scalar-invariant}
 R_E\ge\lambda_r+b_r\sum_{j\in J_r}\Pi_j\quad\hbox{on }D_E.
\end{equation}
Here $R_E$ denotes the stabilized scalar expression, and integrals and
Laplacians use the product of the old warping functions as weight. The
finiteness statement below makes these meanings unambiguous. A nonpositive
counting threshold means the whole component.

\begin{definition}[Finite exceptional data]\label{countable:finite-data}
At stage $r\ge1$, on each whole component $E$ there are finite sets
$I_E\subset J_r$ of label indices and $F_E$ of old warping function indices such
that, for $j\notin I_E$,
\begin{equation}\label{countable:constant-data}
 \mu_j\equiv-\infty,\qquad U_j^-=E,\qquad U_j^+=\varnothing,
 \qquad \Pi_j\equiv c_rp_j^*.
\end{equation}
Every old warping function indexed outside $F_E$ is identically one.
The number $c_r$ is the same on all components and is prescribed by
\begin{equation}\label{countable:constant-data-multiplier}
 c_0=1,\qquad c_{r+1}=(1+3\eta_r/2)c_r.
\end{equation}
For any finite diagram used in a comparison, take the union of its
finite exceptional sets, adding the current label explicitly. This
one finite union then works simultaneously on all its whole components.
At stage zero the analogous property is required on every compact
radial block and buffer, with $c_0=1$.
\end{definition}

The sets $I_E$ need not consist exactly of the nonconstant entries: an
entry may be retained in the exceptional list even if it happens to be
constant. This convention avoids any need to discard previous exceptions.
Only the existence of a finite list matters.

\begin{proposition}[The countable compatible construction]\label{countable:construction}
For the initialized data above, there is a successive compatible
split along regular zero sets for every label. At every finite stage the
component diagram has the invariant above and no nonempty component of
dimension at most one. After the radial stage it has finite exceptional
data on whole components; at stage zero this assertion holds on buffered
compact radial blocks.
The finite reconstruction $X_r$, obtained by truncating the product
ends and gluing the corresponding boundary faces selected by stars, has a proper collapse $q_r:X_r\to M$
and closed signed sides $P_i^{(r)},N_i^{(r)}$ for each processed center
label, with
\[
 P_i^{(r)}\cup N_i^{(r)}=X_r,\qquad
 P_i^{(r)}\cap N_i^{(r)}=W_i^{(r)},
\]
where $W_i^{(r)}$ is its boundary. The forced preimages satisfy
\[
 q_r^{-1}(A_i)\subset P_i^{(r)}\setminus W_i^{(r)},\qquad
 q_r^{-1}(B_i)\subset N_i^{(r)}\setminus W_i^{(r)}.
\]
These inclusions hold in open sides and persist at every subsequent
finite stage. No $n-1$ distinct boundaries, even including the radial boundary,
have a common point. Each center boundary lies over its original open
outer ball, and its components are compact and locally finite.

Old finite truncation levels of the product ends can be
enlarged by identifications preserving collapse,
all boundaries, and all signed indices. For any prescribed sequence
$\eps_r>0$ with $\sum_r\eps_r<1$, the next finite collar compression
can be chosen with collapse displacement less than $\eps_r$, preserving
the previously decided boundary strata and signed sides. These assertions concern each finite stage.
\end{proposition}

\begin{proof}
Initialization is \Cref{countable:initial}. If the construction has
reached stage $r$, its parents have dimension at least two: a first
nonempty output of smaller dimension would have dimension exactly
one, since a zero set decreases dimension by one and a sign preserves it.
\Cref{countable:step} produces the next invariant, and
\Cref{countable:dimension-one} excludes that first output. Thus the
induction continues through every finite stage.

Apply \Cref{component:reconstruction} to each finite initial segment.
Its hypotheses are the compatible splits along regular zero sets and the compact
radial components already established. Its index pieces are exactly
the corresponding compact component truncations. If $n-1$ boundaries met,
assign those labels star entries and choose a closed sign index for
each other processed label containing the point. The corresponding
nonempty truncation would have dimension one, a contradiction.
Forced signs persist under the continuous collapse recursion and
enlargement of the prescribed open sets. In particular
\[
 q_r(W_i^{(r)})\subset M\setminus(A_i\cup B_i)
                         \subset B_g(x_i,2D).
\]
Properness of $q_r$ and compactness of the original closed outer ball
make each boundary component compact; regular local boundary charts give
local finiteness, as in \Cref{component:reconstruction}.

The truncation level identifications in \Cref{component:reconstruction} preserve
collapse, boundaries, and signed indices. Its collar compressions preserve
the old strata and their sides, and \Cref{component:small-displacement}
allows their collapse oscillations to be bounded by the prescribed
$\eps_r$. These choices are made separately on the locally finite
compact component family. Section~\ref{sec:reconstruction} chooses these homeomorphisms
coherently on one fixed copy of $M$, proving \Cref{topology:partitions}.
\end{proof}

\subsection{Initial functions and radial confinement}

\begin{lemma}[Initial functions and radial blocks]\label{countable:initial}
The fixed stretched metric and its initial warping functions admit stage-zero
functions and prescribed open sets satisfying
\eqref{countable:function-invariant}--\eqref{countable:scalar-invariant}
with the prescribed potentials $\Pi_i=p_i^*$. The radial forced signs can be
chosen so that every regular radial split preserving them has locally
finite compact boundary and closed-side components before cylinder
attachment. On every compact radial block and buffer, only finitely
many center labels fail to be identically the negative pole with full
negative prescribed open set, and only finitely many initial warping functions fail to be one.
\end{lemma}

\begin{proof}
For a positive number $p$, the decreasing crossing
\[
 2|y'|=.5y^2+p
\]
from $+\infty$ to $-\infty$ has length
\[
 T(p)=\int_{-\infty}^{\infty}\frac{2\,\dd y}{.5y^2+p}
      =2\pi\sqrt{2/p}.
\]
Extend it by its poles. For center $i$, begin this crossing at stretched
distance $2$ from $A_i$ and compose it with $\dist_h(A_i,\cdot)$.
The distance is $1$-Lipschitz, so the resulting function obeys the
inequality with $.5$ in place of $a_0=.60$. Take prescribed open sets
\[
 U_i^+=\{\dist_h(A_i,\cdot)<1\},\qquad
 U_i^-=\{\dist_h(A_i,\cdot)>T(p_i^*)+3\}.
\]
Their closures lie in the respective pole interiors. The stipulated
separation puts $B_i$ strictly inside the negative prescribed open set. In particular, on
the complement of the original closed ball $\overline B_g(x_i,2D)$,
the function is the negative pole and the negative prescribed open set is full.

The radial label is constructed as in the proof of
\Cref{component:initialization} in Appendix~\ref{sec:components}, with the
potential \(\pi_0\) replaced by \(p_0^*\) and \(T_0\) by \(T(p_0^*)\).
The chartwise smooth-radius construction there, with its
additive constant \(2\) subtracted, gives a smooth proper function
\(\rho:M\to\R\), bounded below, such that
\(|\rho-\dist_h(o,\cdot)|<1\) and \(|\grad\rho|_h<2\).
It uses only local smoothness and completeness, so applies to the
present stretched metric.
The same proof builds alternating positive and negative plateaus of
\(\rho\), separated by transition intervals wider than \(2T(p_0^*)\),
on which the rescaled crossing satisfies the function inequality with
coefficient \(.5\) and potential \(p_0^*\); closures of the prescribed open sets strictly
inside the plateaus; and compact smooth barrier domains, bounded by
regular levels, whose closures stay in the same pole plateaus and
whose value intervals tend to infinity.
The same proof shows that a boundary respecting these signs has compact
components, each lying in one block between consecutive separating
plateaus, that every component of either closed side is trapped between
two separating regions of the opposite sign or lies in the initial compact
block, and that both component families are locally finite. When \(M\)
is compact, it uses a radial label with constant positive pole and full
positive prescribed open set; its split has no boundary and its only side is compact.

The scalar requirement initially holds everywhere because
\[
 \lambda_0+b_0\sum_i p_i^*=.25+1.12/20=.306<.4.
\]
Finally a compact block with a compact buffer meets only finitely many
original closed outer balls and finitely many initial supports of the logarithms of warping functions.
Every other center label and warping function has the asserted constant data there.
\end{proof}

\subsection{Finite exceptional data on entire components}

A \emph{downward diagram} of a component consists of it and all its further
end bases. An \emph{upward diagram} consists of it and all its incident higher-dimensional components.
Each is finite at a fixed stage, by \Cref{component:split}. A finite union
of either type is finite. The constructions below make finitely many
directed passes through such diagrams.

Radial confinement gives the geometric reason to expect finite data on
a whole component. At the first split, a compact radial side or boundary
component is prolonged by cylinders; its entire collapse image stays
in that original compact component. At every later split a child $A$
of $E$ has
\begin{equation}\label{countable:collapse-recursion}
 q_A=q_E\circ p_A,\qquad
 p_A=\begin{cases}
 \text{inclusion},&\text{on the boundary or bulk},\\
 \text{projection to the boundary},&\text{on a fresh cylinder}.
 \end{cases}
\end{equation}
Hence $q_A(A)\subset q_E(E)$, including in all simultaneous normal
directions. A finite incidence diagram therefore has image in a finite
union of compact original radial blocks. Only finitely many original
outer balls and supports of the logarithms of the initial nonunit warping functions meet that image.
Enlargement of the prescribed open sets changes no collapse map, and regularization chooses new
boundaries inside old parents, so both preserve this confinement.

This argument controls the initial supports. Later smooth data need
not factor through $q_A$: they also use data copied from end bases. We
must therefore retain finite exceptional lists through those directed
dependencies, as well as preserve every other entry exactly.

\begin{lemma}[Consequences and directed dependencies]\label{countable:locality}
At a finite stage after the radial split, assume the exact component
geometry of \Cref{component:category} and the finite exceptional data of
\Cref{countable:finite-data}. Then:
\begin{enumerate}[label=(\roman*),leftmargin=*]
\item The sum of potentials on a whole component is a smooth function
with bounded derivatives of every fixed order. The undecided counts,
required sets, and sets $H_E^*$ use only a finite list and are closed.
\item On each whole component and each finite comparison diagram, the old
weight, stabilized scalar, stability forms, and weighted Laplacian are
computed with one finite list of warping functions, after padding with unit factors.
The conclusions of a positive injectivity radius, uniform bounds on
curvature and all its covariant derivatives, and polynomial weighted growth in
\Cref{component:geometry} apply, including on any finite Euclidean product.
\item Compatible choices with positive componentwise tolerances can be
made in the following order: fix individual requests, propagate numerical
requests down the finite chains through higher-dimensional components, and then construct data upward
from deeper end bases. Enlarging an incident higher-dimensional component's exceptional list by the union
over its downward source diagram suffices for the last pass. No data
are propagated back from a newly constructed incident higher-dimensional component to an already
constructed end base.
\end{enumerate}
\end{lemma}

\begin{proof}
On $E$, the sum of potentials is exactly
\begin{equation}\label{countable:potential-sum}
 \sum_{j\in J_r}\Pi_j
 =\sum_{j\in I_E}\Pi_j+c_r\sum_{j\in J_r\setminus I_E}p_j^*.
\end{equation}
The second term is a finite spatial constant. Indeed
\[
 1\le c_r\le\exp\Bigl(\frac32\sum_{s<r}\eta_s\Bigr)
 \le e^{3/400}.
\]
The first term is a finite sum; its derivative bounds follow from
compatible smoothness and the compact-part/product-end decomposition
in \Cref{component:geometry}. On a finite comparison diagram choose common
product thresholds for its finite exceptional list. Every other entry is already a spatial
constant with the same value on each copy. The full sum is therefore
an exact end base pullback beyond those thresholds. Entries outside the finite exceptional list have empty undecided set. Thus a set requiring at least $a\ge1$ undecided labels
is the finite union, over $a$-element subsets of $I_E$, of the intersections of
their closed $G_j$'s. This proves closedness, including for the sets $H_E^*$
with threshold $m-2$. Threshold zero gives $E$.

For (ii), take the union of the finitely many relevant old warping function
indices on the finite diagram and set every missing entry to one.
Writing $u_\alpha=\log f_\alpha$ and $U=\sum_\alpha u_\alpha$, the
scalar expression is
\[
 R_E=\Scal_{h_E}-2\Delta_{h_E}U-|\dd U|_{h_E}^2
                        -\sum_\alpha|\dd u_\alpha|_{h_E}^2.
\]
A padded logarithm is zero; the product weight is unchanged. The
weighted Laplacian is
$\Delta_{h_E}+\langle\dd U,\dd\,\cdot\,\rangle_{h_E}$ and is
unchanged as well. Weighted perimeter and its first and second
variations depend on the same weight. Equivalently, a flat unit
factor adds zero second fundamental form and zero curvature to the
invariant variation calculation. Thus the boundary and interface forms also
agree. The finite-diagram geometric lemma now gives completeness,
a positive injectivity radius and uniform bounds on curvature and all its covariant derivatives, positive lower and finite upper bounds for the
weight, and center-uniform polynomial volume growth. It gives the
same statements after taking a finite Euclidean product. The
regularity problem for a minimizing Caccioppoli set remains on the base, with
singular dimension at most $m-8$ by \Cref{cut:gmt}; padding does not
alter that dimension.

For (iii), a particular end base receives requests from only its finitely
many incident higher-dimensional components and chains through higher-dimensional components. Allowing for the finitely many fixed
product and cutoff bounds, the minimum of the requested positive
errors is positive. Settle all these numbers before constructing the
approximations. A request only decreases an error tolerance: it does
not copy any incident higher-dimensional component value into the end base. During the subsequent upward
pass, the inputs on a component are its own data and already constructed
copies from its finite downward diagram. Their exceptional sets have
a finite union. Outside that union every copy in
\eqref{countable:constant-data} agrees with every other copy. A newly
constructed incident higher-dimensional component is not used to reconstruct its sources. The same
argument applies to finitely many successive directed passes, and
proves the stated dependency rule.
\end{proof}

\subsection{One countable invariant step}

\begin{proposition}[Preservation of the invariant]\label{countable:step}
Suppose that all nonempty parents at stage $r$ have base dimension at
least two and satisfy \eqref{countable:function-invariant}--\eqref{countable:scalar-invariant}.
At $r=0$ use the compact radial blocks of \Cref{countable:initial}; at
$r\ge1$ suppose also that the component geometry and
\Cref{countable:finite-data} hold. There is a compatible split along regular zero sets of label $r$ preserving the
forced signs and satisfying
\eqref{countable:function-invariant}--\eqref{countable:scalar-invariant},
the component geometry, and \Cref{countable:finite-data} at stage $r+1$. Exactly $399$ new warping function factors, with
common indices, are appended at this stage. Each whole resulting
component has a finite relevant old warping function list.
\end{proposition}

\begin{proof}
The finite exceptional lists let each geometric operation use only
finitely many nonconstant functions on a whole component. The remaining
functions and potentials pass through the operation as their prescribed
constants. We first choose the current cut using its own positive floor,
then carry these two types of data through enlargement of the prescribed open sets and regularization.

\emph{The current prescribed weighted mean curvature and minimizing Caccioppoli sets.}
The prescribed weighted mean curvature construction of \Cref{cut:prescribed-curvature} uses a positive floor
only for the label being processed. Replace its finite-family floor
by $p_r^*$. Compatible clipping heights can first be chosen above the
local barrier bounds, with a uniform extra barrier margin of one.
Then use local smoothing errors to obtain, for the smoothed prescribed weighted mean curvature
$p_0$ of \Cref{cut:prescribed-curvature} (its subscript does not refer to the radial
label $0$),
\[
 2|\grad p_0|\le .65p_0^2+1.01\Pi_r.
\]
All smoothing requests involve only this one function and finite
incidence diagrams. The elementary inequality
\[
 .67(y+c)^2-.65y^2\ge-21.775c^2
\]
shows that one common nonempty shift interval, for example
\begin{equation}\label{countable:shift}
 |c|<\min\{.01\sqrt{p_r^*},1/2\},
\end{equation}
preserves both $2|\grad(p_0+c)|\le .67(p_0+c)^2+1.02\Pi_r$ and
the barrier inequalities. Floors for other labels do not enter this
interval.

After the radial stage, \Cref{countable:locality} supplies the positive injectivity radius, uniform bounds on curvature and
all its covariant derivatives, and weighted polynomial-growth hypotheses of
\Cref{cut:generic-uniqueness}. There are countably many components,
Euclidean product dimensions, and tests in that lemma's exhaustion.
Remove the countable exceptional shift sets and choose a single $c$
in \eqref{countable:shift}. The resulting minimizers have the
common product limits of \Cref{cut:end-convergence}; choose them before
any enlargement of the prescribed open sets. The smooth global
positive weight and locally bounded smooth prescribed weighted mean curvature give local BV
existence and strict compactness by \Cref{cut:bv-compactness}. The
relatively compact radial barriers force the Caccioppoli sets by compact
comparisons. All variations and tests lie in buffered compact regions
with finite warping function lists, so the formulas in \Cref{cut:raw-domination}
apply. The GMT and singular-set dimensions are the base dimensions at
this first step as well.

\emph{Smooth approximations, sets $H_E^*$, and enlargement of the prescribed open sets.}
On each post-radial component retain its finite exceptional list, enlarged
by its finite downward source diagram and by the current label.
Labels outside it are already forced everywhere and require no shell
addition. Choose compatible smooth functions $h_j$ within $.1$ of
$\max(-1,\min(1,\mu_j))$ as in \Cref{local:smooth-approximations}, retaining the exact
value $-1$ when $\mu_j=-\infty$ identically. For each component of dimension $m\ge8$ and each set of $m-7$
remaining labels, the image of its singular set under the corresponding
map of smooth approximations has $(m-7)$-dimensional Lebesgue measure zero.
There are countably many such conditions, even if all finite subsets
of the countable label set are included. In the countable product of
uniform probability intervals $(-.3,.3)$, each bad cylinder has
measure zero by its finite-dimensional marginal. A threshold vector $(t_j)$
outside their union therefore exists. At most $m-8$ approximations $h_j$ equal
their thresholds at any singular point. The sign exclusions at poles
follow from the $.1$ approximation to the truncated function, exactly as in
\Cref{local:smooth-approximations}.

Process the finitely many end base depths from deepest upward, keeping the
boundary supports fixed. Suppose every strict end base $B$ of $E$ has $H_B^*$ disjoint from its singular set.
The remaining intersection of its singular set with $H_E^*$ is
\[
 K_E=\Sigma_E\cap
 \{x:\#\{j\in J_r\setminus\{r\}:x\in G'_j\}\ge m-2\}
\]
is closed. It is compact: an escaping sequence would, by
\Cref{component:escape}, have a product limit over a further end base with
compact convergent base coordinates. The finite exceptional list and
closed complements of the prescribed open sets pass the undecided count to that limit.
The end base $B$ has a smaller dimension, so the limit lies in $H_B^*$.
The product convergence of boundary supports forces the limiting point to be
singular, since regular-point graphical convergence excludes nearby
singular points. This contradicts the completed enlargement of the prescribed open sets of that
end base. This argument also treats escape in several normal coordinates.

Apply \Cref{local:plane-containment-pieces,local:shell,local:sharp} on this compact
set $K_E$, only for the finite list which can still be undecided.
Enumerate the signed compact subsets $F$ with the following property:
in Euclidean chart coordinates, for some $r_F>0$, each $F\cap B_r(x)$,
$x\in F$, $0<r<r_F$, lies within $\eta r$ of an affine $(m-8)$-plane,
with $\eta$ as small as required. Choose each shell scale by
\Cref{enlargement:maximum}, hold it fixed, and then send its incident higher-dimensional component starting levels
far enough out by \Cref{enlargement:extension}. In the normal-coordinate function construction
of \Cref{enlargement:normal-functions}, the crossing for label $j$ satisfies
\[
 2|w'|\le .58w^2+p_j^*.
\]
Its length is finite; the inequality is paid by $\Pi_j\ge p_j^*$,
and the rest of that lemma applies unchanged.

Only finitely many additions are required on $E$. Indeed, every point
of the initial compact $K_E$ eventually belongs to a prescribed open set
for each of its strict inequalities $h_j>t_j$ or $h_j<t_j$ among the finite exceptional list.
At most $m-8<m-2$ labels can then remain undecided there. The remaining
subsets of the initial $K_E$ are nested and closed, and their
intersection is empty. Compactness gives a finite index at which the
residual is empty. Stop then. Each specified incident higher-dimensional component receives additions
from finitely many source end bases through finitely many chains. Thus
only finitely many additions and finite minima or maxima affect any
whole component.

The local shell estimate \eqref{local:sharp-estimate} charges
$1.16D_\varepsilon$ at gradient coefficient $.78$. The factor $1.16$
is chosen to exceed $(1+3\eta_r)b_{r+1}$, the coefficient with which
the additions $D_j$ enter the next potentials in the later scalar-form
transfer; see \eqref{regular:eta}.
To control all the countably many components at once, enumerate the pairs
consisting of a component and a possible local addition number by distinct
integers $\nu\ge1$. As in the proof of \Cref{enlargement:global}, assign
$\gamma_\nu=.1\,2^{-\nu}$ to that pair, allocating at most half to
its shell-scale choice and half to its simultaneous finite incident higher-dimensional component
extensions. Starting the bookkeeping form bound at $.78$, the maximum
and extension estimates give
$\theta_*=.78+\sum_{\nu\ge1}\gamma_\nu=.88<.9$ and hence
\begin{equation}\label{countable:enlargement-form}
 \int1.16\sum_{j\ne r}D_jv^2
 \le\theta_*\bigl(\mathcal N(v)+e_r\|v\|_2^2\bigr).
\end{equation}
This is the bound \eqref{enlargement:form-output} of \Cref{enlargement:global}.
Here $\mathcal N(v)$ is the sum of weighted integrals of
$|\nabla v|^2$, with coefficient $1.7$ on a boundary, or coefficient
$.8$ on the bulk and $1.7$ on its attached cylinder. The measure on
each piece is its base volume times the product of the retained
warping functions; $\|v\|_2^2$ is the sum of its weighted squared norms.
Tests are compactly supported away from singularities, and their
two traces match at a joined interface, as in \Cref{cut:reserve}. These are the
same coefficients on every component and product. Only a finite list of
additions occurs on each post-radial component, so the bound persists
after a whole depth has been treated.

Initially, instead work in separated compact transition shells with
compact buffers between radial barriers. Each has finitely many
exceptional labels and warping functions. Its intersection of the singular set with $H_E^*$ is compact, and the same
finite termination applies; there is no incident higher-dimensional component extension yet. Choose
supports within the shell buffers. Their local finiteness makes all
updates locally finite on $M$, and a compact test meets only finitely
many of them. The summable allowances prove
\eqref{countable:enlargement-form} for every such test. Every compact radial
side block meets only finitely many shell buffers. This supplies the
whole-component finite lists after the first attachment.

\emph{Counts and the quantitative transfer.}
Enlargement of the prescribed open sets gives $G'_j\subset G_j$. Only the finite exceptional
list can contribute to the counts on a post-radial component; the analogous
statement holds on buffered compact blocks at the radial step.
Thus \Cref{component:counts} places every new required set in both the old
required set and the set $H_E^*$ disjoint from the singular set under projection, including along
attached cylinders. The finite lists and closed complements of the prescribed open sets also
pass these count inclusions to all further-end base product limits.

The normal tube near $H_E^*$ and comparison constructions of
\Cref{regular:count-sets,regular:partition} use these closed sets $H_E^*$,
regular graphical convergence on product limits, and finite
componentwise derivative bounds. All are now available. Numerical
errors are allocated down chains through higher-dimensional components before the defining functions
are constructed upward.

Apply the quantitative transfer of \Cref{regular:transfer} with the
current $t_r$, $\Delta_r=t_{r+1}-t_r>0$, and $e_0=e_r=\Delta_r/800$.
Its stage margins \eqref{regular:eta} were proved for these general
adjacent parameters.
The potential calculation also applies: on each finite comparison
diagram, \eqref{countable:potential-sum} expresses the old potential
sum as finitely many smooth functions plus one spatial constant.
Only finitely many potential additions near singularities occur there. Thus, with
$\widetilde\Pi_j=\Pi_j+D_j$, the same raw form and comparison estimates
control
\begin{equation}\label{countable:raw-target}
 \lambda_{r+1}+(1+3\eta_r)b_{r+1}
                         \sum_{j\ne r}\widetilde\Pi_j
\end{equation}
near the new required set, with the positive mass and gradient reserves
proved there. The constant tail has no derivative error; potential
approximation errors concern only the finite exceptional list. At the
radial step these estimates are made on buffered compact blocks, as
in that proof.
The interface and potential smoothing can consequently retain an $e_r$
local mass gap. This positive number is fixed for the current stage;
all geometric tolerances may depend on the component.

\emph{Exact constants outside the finite exceptional lists in regularization.}
Restrictions, bulk inclusions, and cylinder projections preserve
identically constant functions. Interface smoothing of an old identically
unit warping function leaves it one. For the remaining potentials, smooth
$(1+3\eta_r/2)\widetilde\Pi_j$ as in
\Cref{regular:functions,regular:constant-preservation}, so that
\begin{equation}\label{countable:smooth-potentials}
 (1+\eta_r)\widetilde\Pi_j\le\Pi_j^{\rm new}
                     \le(1+2\eta_r)\widetilde\Pi_j.
\end{equation}
On a finite exceptional list, allocate the required relative errors
separately, using $p_j^*>0$ for its $j$th entry. If an aggregate error
is prescribed, subdivide it among that finite list. If an entry lies outside every finite exceptional list in the source
diagram, set its output \emph{exactly} to
$c_{r+1}p_j^*$. This is compatible with the smoothing procedure, which
reproduces constants exactly by \Cref{regular:constant-preservation}.

Before blending on an incident higher-dimensional component, enlarge its exceptional list by the union
over its entire finite downward source diagram. Thus an entry on an end base's exceptional list is recorded even if it
was outside the incident higher-dimensional component's own exceptional list.
The remaining copies agree identically. Higher-dimensional component requests merely adjust
numeric tolerances; they do not introduce new source values on end bases.
This proves exact reproduction of the specified constant data on the completed diagram and
precludes a cascade through infinitely many same-depth neighbors.
At the radial step choose a compact buffer for each compact radial
block, enlarged to contain the finitely many chart closures and shell
supports used by its approximations. Only finitely many original
outer balls and supports of the logarithms of the initial nonunit warping functions meet this buffer.
Enlarge the exceptional list by precisely these indices. Every other
entry is the same constant on every averaging neighborhood and is
therefore reproduced exactly before cylinder attachment.

The restriction argument of \Cref{regular:functions}, with the first
inequality in \eqref{regular:eta} and \eqref{countable:smooth-potentials},
gives the next function inequality.
Continuous pullback and enlargement of the prescribed open sets preserve all pole-neighborhood
and forced-preimage conditions.

\emph{Corrections, supersolutions, and the new warping function group.}
The sum in \eqref{countable:potential-sum} is now smooth, bounded, and
exactly compatible, as are its finite-order derivatives. On every
required finite diagram, pad the old warping functions by units. In a product end
over $B$ the operator of \Cref{regular:correction} therefore satisfies
exactly
\[
 L_A=L_B-C\Delta_{\rm normal},\qquad
 L_A=-C\Delta_{\rm old}+R_{\rm old}-V_A^*+K_A,
 \qquad
 V_A^*=\lambda_{r+1}+b_{r+1}\sum_{j\ne r}\Pi_j^{\rm new}.
\]
In particular the potential formula does not depend on the base
dimension or on how many unit factors were padded. The correction is
zero near the new required set and produces a positive form gap on
each whole component. The proof of \Cref{regular:correction} uses finite
minima of the gaps on the component's end base diagram. Its bounded solution of the elliptic equation
and the compatible positive supersolutions in
\Cref{regular:bounded-solutions,regular:solutions} apply with the positive injectivity radius, uniform bounds on curvature and all
its covariant derivatives, and polynomial weighted volume already proved
on each component.
Their tolerances are chosen after the operators and their individual
exact solutions are fixed, and before the final upward blending pass.
Thus their finite-dependency hypothesis is precisely (iii) of
\Cref{countable:locality}.

Let $v_A>0$ be the resulting supersolution with $L_Av_A\ge1/2$.
Append the common group of $k=399$ warping functions
$f=v_A^{2/(k+1)}=v_A^{1/200}$. The scalar identity in
\Cref{regular:step} gives
\[
 R_{\rm new}=R_{\rm old}-C\frac{\Delta_{\rm old}v_A}{v_A}
           =V_A^*-K_A+\frac{L_Av_A}{v_A}
           \ge V_A^*\quad\hbox{on }D_A.
\]
Only the number of \emph{new} factors enters the coefficient $C$.
The old finite list may vary from component to component; unit padding has
no effect. Each stage adds just this one finite group, with the same
factor indices for every component index. All its warping functions are positive and bounded
above and away from zero on a whole component, and exactly compatible.
The initial lists and their finite propagated exceptions remain finite,
so the complete old list at stage $r+1$ is still finite on each whole
component and each finite comparison diagram.

Finally the split along regular zero sets preserves compact truncations and finite
incidence by \Cref{component:split}. Initially the radial separators make
all components compact before prolongation. Compatible smooth metrics
and finite relevant warping functions then give the complete component geometry by
\Cref{component:geometry}. The recursion
\eqref{countable:collapse-recursion} preserves confinement in original
radial blocks. All exceptional-list enlargements just described are
finite, and all constants outside the exceptional lists have multiplier $c_{r+1}$.
This completes the induction step, including the first radial step.
\end{proof}

\subsection{The obstruction in dimension one}

\begin{lemma}[The one-dimensional obstruction]\label{countable:dimension-one}
At any finite stage of the construction, the scalar inequality
\eqref{countable:scalar-invariant} excludes a nonempty one-dimensional
component with the complete geometry and finite warping function list stipulated above.
\end{lemma}

\begin{proof}
Such a component is complete by \Cref{component:geometry} and has only
finitely many relevant positive smooth warping functions $f_1,\ldots,f_q$ on
its whole base. Its arclength universal cover is $\R$, including
when the component itself is a circle. On this line put
$u_\alpha=\log f_\alpha$ and $P=\sum_{\alpha=1}^q u_\alpha$.
The warped-product scalar formula reduces to
\[
 R=-2P''-(P')^2-\sum_{\alpha=1}^q(u_\alpha')^2.
\]
The required count in dimension one is zero, so the invariant holds
everywhere and gives $R\ge\lambda_r>0$. Consequently
\[
 -2P''\ge\lambda_r+(P')^2,\qquad
 \frac{\dd}{\dd s}\arctan\frac{P'}{\sqrt{\lambda_r}}
 \le-\frac{\sqrt{\lambda_r}}2.
\]
Integration over a line interval longer than
$2\pi/\sqrt{\lambda_r}$ contradicts the range of the arctangent.
If $q=0$, the scalar formula is simply $R=0$. The argument takes
place at one finite stage and uses an ordinary finite sum.
\end{proof}


